% Hardware-connection revision of the coordinator's six-page formatted task sheet.
% Existing P/PI starting choices, worked example and targets are retained.
% Build twice with pdflatex. No firmware, controller or website application changes.
\documentclass[11pt,a4paper]{article}
\usepackage[T1]{fontenc}
\usepackage[utf8]{inputenc}
\usepackage{mathpazo,amsmath,amssymb,booktabs,array,tabularx,enumitem}
\usepackage[table]{xcolor}
\usepackage[a4paper,left=18mm,right=18mm,top=20mm,bottom=19mm,headheight=13pt,headsep=9pt,footskip=25pt]{geometry}
\usepackage{fancyhdr,lastpage}
\usepackage[hidelinks]{hyperref}
\hypersetup{pdftitle={EE3600 / EE5610 Practical 2: DC Motor Speed Control with P and PI},pdfauthor={EE3600 / EE5610},pdfsubject={Task sheet with Practical 2 motor-driver wiring and startup}}
\definecolor{ink}{HTML}{233640}
\definecolor{teal}{HTML}{14696C}
\definecolor{muted}{HTML}{5A6870}
\definecolor{pale}{HTML}{EFF6F7}
\definecolor{warning}{HTML}{FFF6EA}
\color{ink}
\setlength{\parindent}{0pt}
\setlength{\parskip}{5pt}
\setlength{\emergencystretch}{2em}
\setlist[enumerate]{leftmargin=1.55em,itemsep=4pt,topsep=4pt,parsep=0pt}
\setlist[itemize]{leftmargin=1.4em,itemsep=3pt,topsep=3pt,parsep=0pt}
\setlength{\abovedisplayskip}{7pt}
\setlength{\belowdisplayskip}{7pt}
\renewcommand{\arraystretch}{1.18}
\pagestyle{fancy}
\fancyhf{}
\fancyhead[L]{\fontsize{9}{11}\selectfont\color{muted}EE3600 / EE5610\quad Automation and Control Systems}
\fancyhead[R]{\fontsize{9}{11}\selectfont\color{muted}Practical 2}
\renewcommand{\headrulewidth}{0.35pt}
\fancyfoot[L]{\fontsize{8.5}{10}\selectfont\color{muted}P and PI speed control\enspace |\enspace Hardware update: 27 October 2026}
\fancyfoot[R]{\fontsize{8.5}{10}\selectfont\color{muted}\thepage\,/\,\pageref*{LastPage}}
\newcommand{\sect}[1]{\vspace{4pt}{\large\bfseries\color{teal}#1}\par\vspace{3pt}}
\newcommand{\sub}[1]{\vspace{4pt}\textbf{#1}\par}
\newcommand{\ui}[1]{\textbf{#1}}
\newcommand{\note}[1]{\par\begingroup\setlength{\fboxsep}{7pt}\noindent\colorbox{pale}{\parbox{\dimexpr\linewidth-14pt\relax}{#1}}\endgroup\par}
\newcommand{\warn}[1]{\par\begingroup\setlength{\fboxsep}{7pt}\noindent\colorbox{warning}{\parbox{\dimexpr\linewidth-14pt\relax}{#1}}\endgroup\par}
\newcolumntype{Y}{>{\raggedright\arraybackslash}X}
\newcommand{\blank}{\rule{27mm}{0.3pt}}
\begin{document}
\fontsize{10.5}{13}\selectfont

% PAGE 1 -- original model handover, with a pointer to the new wiring pages.
{\LARGE\bfseries Practical 2}\par
{\Large\bfseries DC Motor Speed Control with P and PI}\par
Name / group: \rule{93mm}{0.3pt}\hfill Date: \rule{29mm}{0.3pt}
\note{\textbf{Main learning task:} explain why proportional control leaves a speed offset and how integral action removes it. Reuse the Practical 1 model, simulate P and PI, test a reviewed PI candidate, then verify three fresh repeats. Position is optional in simulation only.}
Website: \href{https://ericwangjcu.github.io/EE3600-Motor-Practicals/practical2-control.html}{\textbf{Open Practical 2 -- speed control with P and PI}}\\
{\small\url{https://ericwangjcu.github.io/EE3600-Motor-Practicals/practical2-control.html}}

\sect{1. Carry forward the Practical 1 model}
Select the same motor used in Practical 1. The two website choices are:
\begin{tabularx}{\linewidth}{@{}Y p{29mm}p{37mm}@{}}
\toprule
\textbf{Motor selection}&\textbf{Gearbox ratio, $N$}&\textbf{Motor encoder counts/rev}\\\midrule
6 V / 280 RPM output & 21.3:1 & 44\\
6 V / 1360 RPM output & 4.4:1 & 44\\\bottomrule
\end{tabularx}
Both use the same P2 UF2. The selector changes the gearbox conversion, not your measured model or controller gains.

Use your existing experimental gain $\overline K$, time constant $\overline\tau$ and gearbox ratio $N$. Do not repeat identification, voltage/speed measurements, physical-parameter derivations or P1 model validation.
\[
G_{P1}(s)=\frac{\overline K}{\overline\tau s+1}
\qquad\text{[gearbox-output rad/s per volt].}
\]
The website uses \textbf{motor-shaft RPM} and \textbf{fractional duty} $u$ ($0.40=40\%$). With the actual declared supply $V_s$ and nominal approximation $V_{\mathrm{motor}}\approx V_su$:
\[
G_n(s)=\frac{K}{\tau s+1},\qquad K=\overline K N\frac{60}{2\pi}V_s,\qquad \tau=\overline\tau.
\]
\centerline{\small $K$: motor RPM per duty fraction. $\tau$: seconds.}
The conversion changes the gain, not the time constant. $K$ is the plant gain, not controller gain $K_p$. Keep your own motor/gearbox/load and record the supply and units. Driver losses and nonlinear behaviour are not separately represented by this voltage scaling.
\begin{tabularx}{\linewidth}{@{}YYYYY@{}}
\toprule
$\overline K$ [rad/s/V]&$\tau$ [s]&$N$&$V_s$ [V]&Converted $K$\\[16pt]\bottomrule
\end{tabularx}
\note{\textbf{Website section 1:} select your motor and actual supply, enter your P1 gain and time constant, then choose the agreed target in motor RPM. Select \ui{Use my P1 model}. The model is converted once; no controller gains are applied.}
Follow the website in order: 1 P1 model $\to$ 2 P only $\to$ 3 Add I $\to$ 4 Targets $\to$ 5 Pico test $\to$ 6 Results. Use \ui{Next/Back} or the numbered sections.
\textbf{Before the Pico test, complete the new hardware and startup instructions in Sections 5A--5B. P2 motor-power wiring is different from P1.}

\newpage
% PAGE 2 -- preserve the Codex P starting choice.
\sect{2. Analyse and simulate P-only control}
Use unity negative feedback, $e=r-n$, with $C(s)=K_p$. Derive the closed-loop response:
\[
T_P(s)=\frac{KK_p}{\tau s+1+KK_p},\qquad \tau_{\mathrm{cl}}=\frac{\tau}{1+KK_p},
\]
\[
n_\infty=\frac{rKK_p}{1+KK_p},\qquad \frac{e_{\mathrm{ss}}}{r}=\frac{1}{1+KK_p}.
\]
\sub{Choose an initial P gain}
Use the converted plant gain $K$ from section 1. Choose a modest initial loop gain:
\[
KK_{p,0}=0.5,\qquad K_{p,0}=\frac{0.5}{K},\qquad n_\infty=\frac r3,\qquad \tau_{\mathrm{cl}}=\frac{\tau}{1.5}.
\]
This deliberately demonstrates offset. The factor 0.5 is a teaching starting choice, not an optimal-gain formula. Gain units are duty/RPM.
\begin{enumerate}
\item Set $K_i=K_d=0$. Enter $K_{p,0}$ as the base P gain. Compare $0.5\times$, $1\times$ and $2\times$ this base at the same reference.
\item Plot reference and speed; record final speed, offset, time constant and required duty.
\item Explain why P needs a remaining error to supply the drive needed for steady rotation.
\end{enumerate}
For this ideal first-order model, positive P gives a monotonic response with \textbf{no overshoot}. Increasing P speeds the response and reduces, but does not eliminate, the offset. Settling around the wrong final speed is not successful target tracking. \textbf{A physical P-only trial is not required.}

The control website's \ui{Compare P} panel uses the ideal model. The separate legacy limited simulator includes a fixed actuator lag, so its P response need not have the same no-overshoot property.

\newpage
% PAGE 3 -- preserve the Codex integral-time starting choice and worked example.
\sect{3. Add I and tune PI in simulation}
\[
C_{PI}(s)=K_p+\frac{K_i}{s},\qquad
T_{PI}(s)=\frac{K(K_ps+K_i)}{\tau s^2+(1+KK_p)s+KK_i}.
\]
I accumulates persistent error. At zero error, the accumulated integral contribution can remain nonzero and supply the required steady drive. Zero ideal step error assumes a stable loop and an achievable target.
\sub{Choose an initial I gain, then refine}
Keep your selected P gain fixed and D = 0. Express PI using the integral time:
\[
C(s)=K_p\left(1+\frac{1}{T_is}\right),\qquad K_i=\frac{K_p}{T_i}.
\]
$T_i$ is the integral time in seconds, not the settling time. For constant error, the integral contribution equals the proportional contribution after this time.
\[
T_{i,0}=5\tau,\qquad K_{i,0}=\frac{K_p}{5\tau}\quad\text{[duty/(RPM$\cdot$s)].}
\]
The factor 5 starts integral action slowly relative to the plant. It is a simulation heuristic, not a lecture-derived optimum or an approved hardware setting.
\begin{enumerate}
\item Enter $K_{i,0}$ as the base I gain and compare $0.5\times$, $1\times$ and $2\times$ it, holding P fixed. Inspect speed, duty, overshoot, rise, settling and final error.
\item If tracking is too slow, increase I in simulation. If overshoot or oscillation grows, reduce I. Check the whole trace; a smaller final error alone is not a pass.
\item If I changes alone cannot meet the targets, adjust P one step at a time and repeat the I comparison. Record each candidate, its results and the reason for the change; save the initial and selected designs.
\item Check the complete response, including the PI zero at $-K_i/K_p$. Denominator-only overshoot formulas do not determine this response.
\item Before hardware, check duty limits, the 10 ms controller interval and declared filtering in MATLAB/Simulink. Do not add measurement lag twice if included in the P1 fit. Obtain tutor review before entering hardware gains.
\end{enumerate}
\textbf{Worked simulation start only:} if $K=5955.553$ and $\tau=0.051$ s, $K_{p,0}=0.00008396$, $T_{i,0}=0.255$ s and $K_{i,0}=0.0003292$. Use your own $K$ and $\tau$. If you select a different P, recalculate the starting I from that P. These are not the final bench gains.
\note{\textbf{PI is sufficient when it meets the requirements. D was not needed in the reported successful speed tests.} No extra time-constant helper or pole-choice exercise is required.}

\newpage
% PAGE 4 -- same assessment conventions.
\sub{How the three lecture methods apply}
\begin{tabularx}{\linewidth}{@{}p{39mm}Y@{}}
\toprule\textbf{Method}&\textbf{Use in this practical}\\\midrule
Model-based analysis&Derive P/PI behaviour and explain gain effects.\\
Computational tuning&Adjust or search gains in simulation. State the settings and criteria, then validate one reviewed candidate.\\
Ziegler--Nichols&Explain both methods and assess applicability. The ideal model has no positive reaction-curve delay and no finite positive ultimate gain. Use a separate suitable tutor example for numerical practice.\\\bottomrule
\end{tabularx}
\textbf{These are not three compulsory hardware exercises.} Do not invent a delay or ultimate gain, or increase physical gain to force sustained oscillation.
\sect{4. Speed requirements and measurement conventions}
Use one agreed achievable positive step from rest, with the same target, supply, load and limits throughout. The reported reference kit used 1000 motor RPM at 6 V; this is not a universal target for every kit.
\begin{tabularx}{\linewidth}{@{}p{32mm}Y@{}}
\toprule\textbf{Quantity}&\textbf{Requirement}\\\midrule
Overshoot&$\leq5\%$ of the commanded step, beyond the final target.\\
10--90\% rise&$\leq0.50$ s. Also report command-to-10\% time to expose startup delay.\\
Settling&$\leq1.00$ s from the command event; all later assessment samples in band, with $\geq2$ s final in-band dwell.\\
Final error&Last-2-second mean absolute error $\leq\max(2\%\text{ of target},2q)$. Report signed bias too.\\\bottomrule
\end{tabularx}
\sub{Encoder-aware assessment band}
For $C_m$ decoded counts per motor revolution and the 50 ms speed-measurement window:
\[
q=\frac{60}{0.050C_m}\quad\text{motor RPM}.
\]
The settling band is centred on the commanded final target $r_f$:
\[
r_f\pm\max\bigl(0.02|r_f-r_0|,\,2q\bigr).
\]
For $C_m=44$, $2q=54.545$ motor RPM. At a 1000 RPM target this is $\pm5.45\%$, \textbf{not strict $\pm2\%$}. State which term governs the band. The ideal simulation comparison reports strict 2\% separately.
\note{\textbf{Assess against the commanded target.} Measure time from the command, not from the start of motion. Faults, incomplete recordings or insufficient final dwell cannot pass.}
\begin{tabularx}{\linewidth}{@{}YY@{}}
\toprule\textbf{Agreed target: motor RPM}&\textbf{Assessment-band half-width: motor RPM}\\[15pt]\bottomrule
\end{tabularx}

\newpage
% PAGE 5 -- NEW: physical wiring, grounded in P1 wiring, P2 firmware and board documentation.
\sect{5A. Hardware connections -- what changes from P1}
\note{\textbf{Practical 1:} the bench supply drove the motor directly; the Pico measured the encoder.\\
\textbf{Practical 2:} the bench supply powers the \textbf{motor-driver board}, and the motor connects to \textbf{Motor 1}. The Pico now controls the drive as well as measuring speed. The encoder wiring stays the same.}
\sub{Make all connections with power removed}
Switch \textbf{OFF} the bench supply and the driver-board switch, and unplug the Pico USB cable before moving wires. Secure the motor with the shaft clear.
\begin{enumerate}
\item \textbf{Fit/check the Pico.} Use the original RP2040 Pico seated in the Kitronik 5331 board's two header sockets. Match the Pico outline/orientation on the board; check for reversed or offset pins. Do not remove an already-correct fitted Pico.
\item \textbf{Remove the direct P1 motor-supply connection.} Disconnect \textbf{both motor winding leads} from the bench supply. Move the motor's \textbf{red and white leads} to the \textbf{two Motor 1 output terminals}, one lead in each terminal. Use the tutor-marked lead order; leave Motor 2 unused.
\item \textbf{Power the board, not the motor directly.} Connect bench supply \textbf{+ to the board POWER +} terminal and bench supply \textbf{-- to POWER -- / GND}, following the board markings. For the two 6 V kits, use a \textbf{6 V supply} unless the tutor explicitly approves a different configuration. Leave the supply and board switch OFF for now.
\item \textbf{Retain/check the encoder wiring} using the table below. Keep the encoder's low-voltage wiring separate from the motor-power terminals. If wire colours differ, identify the signals with the tutor rather than guessing.
\end{enumerate}
\begin{tabularx}{\linewidth}{@{}p{37mm}p{37mm}Y@{}}
\toprule\textbf{Wire on the class kit}&\textbf{Pico connection}&\textbf{Purpose}\\\midrule
Yellow&GP26&Encoder A signal\\
Green&GP27&Encoder B signal\\
Blue&3V3 (OUT), or board 3V breakout&Encoder supply; not motor supply\\
Black&GND&Encoder ground, common with Pico/board ground\\\bottomrule
\end{tabularx}
{\small GP26 and GP27 are GPIO names, not physical pin numbers. Retain the existing P1 harness; do not substitute GP28 for GP26. The board's header sockets already carry the Motor 1 control signals GP2/GP3 and common ground; no extra motor-control jumpers are required.}
\warn{\textbf{Do not leave the bench supply connected directly to the motor.} Do not connect either Motor 1 output to supply + or GND, and do not power the motor from Pico 3V3 or a GPIO. The encoder must not receive the 6 V motor supply. USB alone is not the motor-power source for this practical.}
\textbf{Connection check:} bench supply $\to$ board POWER input; Motor 1 $\to$ red/white motor leads; encoder $\to$ GP26/GP27, 3V3 and GND. The PC connects to the \textbf{Pico USB port} with a data-capable cable in Section 5B.

{\fontsize{8}{10}\selectfont\color{muted}Connection basis: P1 task sheet, Section 3 (class-kit colours); P2 hardware guide/firmware (Motor 1 and GP26/GP27); \href{https://resources.kitronik.co.uk/pdf/5331-compact-motor-driver-raspberry-pi-pico-datasheet.pdf}{Kitronik 5331 datasheet}, pp. 1--2 (board power, sockets and Motor 1).}

\newpage
% PAGE 6 -- NEW: firmware handover, connection and controlled power-up.
\sect{5B. Load P2 firmware and bring the hardware online}
\sub{1. Replace P1 firmware once, if needed}
The P1 encoder-only firmware does not run this controller. Use \textbf{EE3600\_Practical\_2\_Pico.uf2} for \textbf{both motor versions}. If the Pico already reports \textbf{EE3600\_P2 1.0.0}, skip re-flashing.

Download from the website's \ui{P2 firmware} link, or \href{https://ericwangjcu.github.io/EE3600-Motor-Practicals/firmware/EE3600_Practical_2_Pico.uf2}{\textbf{open the P2 UF2 download}}.
\begin{enumerate}
\item Keep the bench motor supply disconnected/OFF and the driver-board switch OFF. Unplug the Pico USB cable.
\item Hold \textbf{BOOTSEL} on the Pico while reconnecting its USB cable to the PC. Release BOOTSEL after the \textbf{RPI-RP2} drive appears.
\item Copy \textbf{EE3600\_Practical\_2\_Pico.uf2} onto RPI-RP2. The drive disappears and the Pico reboots. Do not copy the P1 recovery UF2 or install MicroPython for this task.
\end{enumerate}
\sub{2. Connect to the website while motor power is OFF}
\begin{enumerate}
\item Open the Practical 2 website in desktop \textbf{Chrome or Edge}. Close Thonny, serial monitors and other pages using this Pico's port.
\item In website section 1 select the \textbf{actual motor} (280 RPM / 21.3:1 or 1360 RPM / 4.4:1) and \textbf{actual supply voltage}. Complete the model and simulation work in sections 1--4.
\item Go to website section 5. Click \ui{Connect Pico}, select the Pico serial port and grant the browser permission. Confirm compatible P2 firmware and a \textbf{DISARMED} state. Connecting must not start the motor.
\item Check \ui{Tutor settings}, including the maximum duty and normal \textbf{10 s trial}. Click \ui{Apply configuration} while disarmed, then check the readback: $C_m=44$, the selected gearbox ratio, actual supply and agreed limits. Do not enter approximately 937 output-shaft counts as motor counts.
\end{enumerate}
\sub{3. Check, then switch on motor power}
\begin{enumerate}
\item Ask the tutor to check the wiring, approved supply-current limit, duty ceiling and motor/encoder direction pairing. The board datasheet specifies \textbf{1.5 A maximum per motor channel}; a supply marked 3 A is not an instruction to set a 3 A current limit. Do not carry out a stall test.
\item With the shaft clear, enable the bench output and the board switch. Check that the Pico still reports \textbf{DISARMED}, with \textbf{zero applied duty and zero speed}. If switching power resets the Pico, reconnect and reapply the configuration before continuing.
\item Enter the \textbf{reviewed PI candidate}, with $K_d=0$. Start only by selecting \ui{Arm \& run speed trial}. Observe speed and applied duty; keep \ui{STOP} accessible.
\end{enumerate}
\warn{For a positive command, motion must give positive speed feedback. If the motor runs but feedback is zero/negative, it runs away, stalls, or a fault appears: select \ui{STOP}, switch off motor power and call the tutor. Do not correct wiring while powered, increase limits, bypass a fault or repeatedly re-arm. STOP removes drive; the shaft can coast.}
\textbf{At the end:} STOP; wait for DISARMED and coast-down to zero speed; switch the board and bench supply OFF; save the recording, then disconnect the website/USB. Do not touch or rewire a moving motor.

\newpage
% PAGE 7 -- original website sequence, now referring to the hardware preflight.
\sect{5C. Website and Pico workflow}
The website sections 1--6 follow this sheet. Only one section is shown at a time; use \ui{Next/Back} or the numbered section buttons. The motor selector is in section 1. \textbf{Complete Sections 5A--5B before the physical test.}
\begin{enumerate}
\item \textbf{Section 1 -- P1 model.} Select your motor and actual supply, enter your own $\overline K$ and $\tau$, and choose the agreed target in motor RPM. Select \ui{Use my P1 model}. The model is converted once for the following simulation sections; no gains are applied.
\item \textbf{Section 2 -- P only.} Enter a base P gain and compare $0.5\times$, $1\times$ and $2\times$ P. Save the simulated CSV and record the offset and speed of response.
\item \textbf{Section 3 -- Add I.} Keep the selected base P fixed; compare three I settings. Refine and justify your candidate. Use MATLAB/Simulink to check the declared limits and sampled behaviour before tutor review.
\item \textbf{Section 4 -- Targets.} Review the response targets and encoder-aware band in section 4 of this sheet.
\item \textbf{Section 5 -- Pico test.} Use desktop Chrome/Edge and the P2 UF2. Select \ui{Connect Pico}, then \ui{Apply configuration} while disarmed. The chosen profile sets $N$ and motor counts per revolution automatically; check the displayed readback and tutor-approved duty limit. Enter your reviewed $K_p$ and $K_i$. D remains zero and the normal trial is 10 s.
\item \textbf{Arm \& run speed trial.} Inspect speed and duty. \ui{STOP} removes drive; allow coast-down. A fault requires tutor diagnosis, not an automatic retry.
\item \textbf{Section 6 -- Results.} Select \ui{Analyse response $\to$ Save CSV}. Save before another trial. Stop before changing a gain; document each refinement.
\item \textbf{Freeze and repeat three times.} Retain every attempt, reload saved P2 CSVs and check reproducible metrics. Do not change gains between final repeats.
\end{enumerate}
\note{\textbf{Gain units are the same in the simplified comparison and hardware panels.}\\
$K_p$: duty fraction/RPM.\quad $K_i$: duty fraction/(RPM$\cdot$s).\\
Only the separate legacy simulator uses percent duty/output RPM: divide those gains by $100\times N$ before hardware entry. Unit conversion alone does not match filtering or sampling.}
\warn{Do not apply a 67\% minimum-duty clamp or hidden starting boost. The report only establishes that 67\% produced movement in that short check, not an exact threshold. A 100\% ceiling was approved for the reported 6 V bench, not every kit. \textbf{On a fault, stop and obtain tutor diagnosis before restarting.} Physical position controls are not part of this practical.}
\textbf{Reviewed PI candidate:}\quad $K_p=$ \blank\quad $K_i=$ \blank\quad $K_d=0$.

\newpage
% PAGE 8 -- original results, optional simulation and validation qualifications.
\sect{6. Results and submission}
\begin{tabularx}{\linewidth}{@{}p{33mm}p{25mm}YYYY@{}}
\toprule
\textbf{Case}&\textbf{Gains}&\textbf{Over-shoot (\%)}&\textbf{Rise (s)}&\textbf{Settling (s)}&\textbf{Final MAE / bias (RPM)}\\\midrule
P simulations&&&&&\\[14pt]\midrule
PI simulations initial/final&&&&&\\[10pt]\midrule
Initial hardware&&&&&\\[14pt]\midrule
Final repeat 1&&&&&\\[14pt]\midrule
Final repeat 2&&&&&\\[14pt]\midrule
Final repeat 3&&&&&\\[14pt]\bottomrule
\end{tabularx}
Submit carried-forward P1 values and the unit conversion, P/PI analysis, simulation/tuning record, new P2 CSVs, speed/reference/duty plots, the comparison table and discussion. \textbf{Do not resubmit P1 identification work.} Preserve all raw timestamps; exclude only explicitly reported pre-command transition frames from assessment, never startup delay or failed responses.

Explain P's offset, how I supplies steady drive, what increasing I changed, model/hardware differences, whether each final repeat passed, the assessment band, and why the two Ziegler--Nichols methods do not directly fit this ideal model.
\sub{Reported bench reference -- not a prescribed student answer}
{\small One unloaded 6 V kit, target 1000 motor RPM, $N=21.3$ and $C_m=44$: $K_p=0.0001284515583$, $K_i=0.0010074632025$, $K_d=0$. Three fresh reported repeats met numerical speed targets. Worst overshoot: 3.655\%; rise: 0.2833 s; settling: 0.962 s; final MAE: 13.465 RPM; band: $\pm54.545$ RPM. These are supplied report values, not new website measurements. A later $N=21.6$ estimate does not rewrite those historical runs.}
\sect{7. Optional position extension -- simulation only}
\[
G_\theta(s)=\frac{6}{N}\frac{G_n(s)}s\qquad\text{[output degrees / duty fraction].}
\]
{\small Use the selected lecture's position exercise in MATLAB/Simulink: compare P and PD/PID, state the controller structure and reference, and assess its optional simulation criteria (overshoot $\leq5\%$, $0.25\leq$ settling time $\leq0.80$ s and zero ideal final error). Do not copy speed gains directly. Both reported physical position candidates faulted; neither is a validated bench demonstration. \textbf{No position hardware trial is required.}}
\sub{Instructor release status}
{\fontsize{8.5}{10.5}\selectfont
The supplied bench report supports numerical speed performance only at the stated operating point. The revised page needs a physical native-browser trial, CSV export/re-import against real records and a separate physical USB-removal check before complete end-to-end validation is claimed. Browser tests with synthetic CSVs do not establish hardware performance.

Basis: coordinator's six-page formatted task sheet, existing P1 wiring/model, selected PID lecture and supplied bench reports. This revision adds Sections 5A--5B and their cross-references. Initial P/PI calculations, worked example and assessment targets are retained. Lecture slides, Practical 1 and firmware are unchanged. Board connection details were checked against the \href{https://resources.kitronik.co.uk/pdf/5331-compact-motor-driver-raspberry-pi-pico-datasheet.pdf}{Kitronik 5331 datasheet} and the \href{https://github.com/ericwangjcu/EE3600-Motor-Practicals/blob/main/docs/practical2-hardware.md}{P2 project hardware guide}.}
\end{document}
